\documentclass[11pt]{article}

\usepackage[margin=1in]{geometry}
\usepackage[T1]{fontenc}
\usepackage{lmodern}
\usepackage{microtype}
\usepackage{amsmath,amssymb,amsthm}
\usepackage{booktabs}
\usepackage{graphicx}
\usepackage[numbers,sort&compress]{natbib}
\usepackage[colorlinks=true,citecolor=blue,linkcolor=black,urlcolor=blue]{hyperref}

\title{Boundary Density Likelihood for Time-Series Event Detection}
\author{Clark Peng \and Tolga Dincer}
\date{}

\newtheorem{proposition}{Proposition}

\begin{document}
\maketitle

\begin{abstract}
Many time-series event-detection systems are trained as per-timestep segmentation models and then post-processed to recover event boundaries.
This draft reframes interval event detection as boundary-density estimation: annotated onsets and offsets define sparse boundary measures, and neural sequence models learn nonnegative boundary rates by maximum likelihood.
The final submission will compare this likelihood objective against legacy mean-squared-error heatmap targets and segmentation baselines under tolerance-based event average precision.
\end{abstract}

\section{Introduction}

This manuscript is being rewritten for an ML workshop audience.
The central claim is no longer that heatmap-style mean-squared-error regression is universally better than segmentation.
Instead, we study whether a likelihood over event-boundary locations is a better training objective for metrics that score localized onsets and offsets.

\section{Related Work}

The rewrite will connect four threads: time-series event detection, temporal boundary localization, heatmap/keypoint regression, and temporal point-process likelihoods.

\section{Problem Setup}

Let $x_{1:T}$ denote a multivariate time series and let each interval event be represented by an onset and offset pair $(a_j,b_j)$.
The method constructs separate boundary targets for onset and offset channels.

\section{Boundary Density Likelihood}

The main method will define smoothed boundary count targets and train nonnegative model rates with a discrete Poisson negative log-likelihood.

\section{Experiments}

The final experiments will rerun the core sleep and seizure benchmarks from scratch, treating the existing figures as disposable placeholders.

\section{Discussion}

This rewrite will emphasize when boundary likelihoods are appropriate, when segmentation remains the right abstraction, and what empirical evidence is still missing.

\bibliographystyle{plainnat}
\bibliography{main}

\end{document}
